% GENERATED FILE. Edit canonical lessons, then run npm run generate:books.
% canonical-source-hash: fnv1a64:1c03953470f1d442
% canonical-lessons: GE-C86-schwein, GE-C86-vogel, GE-C86-maus, GE-C86-fuchs, GE-C86-schlange

\chapter{Pig, Bird, Mouse, Fox, Snake}
\label{ch:ge-pig-bird-mouse-fox-snake}
\hlchaptermodality{86}

\begin{chapteropening}
\textbf{By the end of this chapter:} \emph{I can say a pig, a bird, a mouse, a fox and a snake.}

Complete the last lesson of chapter 86: I can say a pig, a bird, a mouse, a fox and a snake.
\end{chapteropening}

\section[das Schwein]{das Schwein — a pig}

\label{lesson:GE-C86-schwein}



\begin{quote}
Before the new one: say the German for a horse, then the German for a cow.
\end{quote}

\subsection*{You'll want to know: das Schwein}
\textbf{das Schwein} — \textquotedblleft{}a pig\textquotedblright{}.

\begin{cousinweb}
\textbf{Schwein} $\leftarrow$ Old High German \emph{swīn} $\leftarrow$ Proto-Germanic \emph{*swīną}, \textquotedblleft{}swine, pig\textquotedblright{} — inherited, the same word as English \emph{swine}. (EtymDB 2.1)
\end{cousinweb}

\begin{grammarlens}[title={What you've built}]
One more word for this chapter.
\end{grammarlens}

\subsection*{Your turn}
\begin{itemize}
\raggedright
  \item \emph{Say it:} \emph{das Schwein}
  \item \emph{Say it:} \emph{das Schwein}, once more
  \item \emph{Say it:} say \emph{das Schwein}
  \item \emph{Recall:} say the German for a horse, then the German for a cow, then say \emph{das Schwein} again
\end{itemize}

\begin{tcolorbox}[breakable,colback=teal!4,colframe=teal!35!black,title={Before you move on}]
What does das Schwein mean? (\textquotedblleft{}A pig\textquotedblright{}.) Say it once more.
\end{tcolorbox}
\section[der Vogel]{der Vogel — a bird}

\label{lesson:GE-C86-vogel}



\begin{quote}
Before the new one: say the German for a cow, then the German for a pig.
\end{quote}

\subsection*{You'll want to know: der Vogel}
\textbf{der Vogel} — \textquotedblleft{}a bird\textquotedblright{}.

\begin{cousinweb}
\textbf{Vogel} $\leftarrow$ Old High German \emph{fogal} $\leftarrow$ Proto-Germanic \emph{*fuglaz}, \textquotedblleft{}bird\textquotedblright{} — inherited, the same word as English \emph{fowl}. (EtymDB 2.1)
\end{cousinweb}

\begin{grammarlens}[title={What you've built}]
One more word for this chapter.
\end{grammarlens}

\subsection*{Your turn}
\begin{itemize}
\raggedright
  \item \emph{Say it:} \emph{der Vogel}
  \item \emph{Say it:} \emph{der Vogel}, once more
  \item \emph{Say it:} say \emph{der Vogel}
  \item \emph{Recall:} say the German for a cow, then the German for a pig, then say \emph{der Vogel} again
\end{itemize}

\begin{tcolorbox}[breakable,colback=teal!4,colframe=teal!35!black,title={Before you move on}]
What does der Vogel mean? (\textquotedblleft{}A bird\textquotedblright{}.) Say it once more.
\end{tcolorbox}
\section[die Maus]{die Maus — a mouse}

\label{lesson:GE-C86-maus}



\begin{quote}
Before the new one: say the German for a pig, then the German for a bird.
\end{quote}

\subsection*{You'll want to know: die Maus}
\textbf{die Maus} — \textquotedblleft{}a mouse\textquotedblright{}.

\begin{cousinweb}
\textbf{Maus} $\leftarrow$ Old High German \emph{mūs} $\leftarrow$ Proto-Germanic \emph{*mūs} — inherited, the same word as English \emph{mouse}. (EtymDB 2.1)
\end{cousinweb}

\begin{grammarlens}[title={What you've built}]
One more word for this chapter.
\end{grammarlens}

\subsection*{Your turn}
\begin{itemize}
\raggedright
  \item \emph{Say it:} \emph{die Maus}
  \item \emph{Say it:} \emph{die Maus}, once more
  \item \emph{Say it:} say \emph{die Maus}
  \item \emph{Recall:} say the German for a pig, then the German for a bird, then say \emph{die Maus} again
\end{itemize}

\begin{tcolorbox}[breakable,colback=teal!4,colframe=teal!35!black,title={Before you move on}]
What does die Maus mean? (\textquotedblleft{}A mouse\textquotedblright{}.) Say it once more.
\end{tcolorbox}
\section[der Fuchs]{der Fuchs — a fox}

\label{lesson:GE-C86-fuchs}



\begin{quote}
Before the new one: say the German for a bird, then the German for a mouse.
\end{quote}

\subsection*{You'll want to know: der Fuchs}
\textbf{der Fuchs} — \textquotedblleft{}a fox\textquotedblright{}.

\begin{cousinweb}
\textbf{Fuchs} $\leftarrow$ Old High German \emph{fuhs} $\leftarrow$ Proto-Germanic \emph{*fuhsaz} — inherited, the same word as English \emph{fox}. (EtymDB 2.1)
\end{cousinweb}

\begin{grammarlens}[title={What you've built}]
One more word for this chapter.
\end{grammarlens}

\subsection*{Your turn}
\begin{itemize}
\raggedright
  \item \emph{Say it:} \emph{der Fuchs}
  \item \emph{Say it:} \emph{der Fuchs}, once more
  \item \emph{Say it:} say \emph{der Fuchs}
  \item \emph{Recall:} say the German for a bird, then the German for a mouse, then say \emph{der Fuchs} again
\end{itemize}

\begin{tcolorbox}[breakable,colback=teal!4,colframe=teal!35!black,title={Before you move on}]
What does der Fuchs mean? (\textquotedblleft{}A fox\textquotedblright{}.) Say it once more.
\end{tcolorbox}
\section[die Schlange]{die Schlange — a snake; a queue}

\label{lesson:GE-C86-schlange}



\begin{quote}
Before the new one: say the German for a mouse, then the German for a fox.
\end{quote}

\subsection*{You'll want to know: die Schlange}
\textbf{die Schlange} — \textquotedblleft{}a snake; a queue\textquotedblright{}.

\begin{cousinweb}
\textbf{Schlange} $\leftarrow$ Old High German \emph{slango} $\leftarrow$ Proto-Germanic \emph{*slangô}, \textquotedblleft{}snake\textquotedblright{} — inherited. (EtymDB 2.1) A queue is a line that winds like a snake.
\end{cousinweb}

\begin{grammarlens}[title={What you've built}]
One more word for this chapter.
\end{grammarlens}

\subsection*{Your turn}
\begin{itemize}
\raggedright
  \item \emph{Say it:} \emph{die Schlange}
  \item \emph{Say it:} \emph{die Schlange}, once more
  \item \emph{Say it:} say \emph{die Schlange}
  \item \emph{Recall:} say the German for a mouse, then the German for a fox, then say \emph{die Schlange} again
\end{itemize}

\begin{tcolorbox}[breakable,colback=teal!4,colframe=teal!35!black,title={Before you move on}]
What does die Schlange mean? (\textquotedblleft{}A snake; a queue\textquotedblright{}.) Say it once more.
\end{tcolorbox}
